\roadmapSection{1}{RELAY AUTOMATION PANEL}{2--3 Weeks}

{\small\itshape\color{arligray} Uses the salvaged relay/telecom board's 12 relays to control real appliances from an ESP32 and your phone. This is your first complete, physically satisfying build.}

% ══════════════════════════════════════════════════════════════════════════════
\roadmapSubSection{1.1}{Understanding What You Have}{2 Days}

\roadmapHead{Inspect Before You Desolder Anything}
\checkitem{Count and photograph the two rows of relays on the green board (typically 6 + 6 = 12 relays on this style of telecom card)}{}
\checkitem{Read the part number printed on top of one relay with a magnifier or phone macro shot, and search it online to get its datasheet}{This tells you coil voltage (often 5V or 12V), coil resistance, and contact current rating (often 1--10A)}
\checkitem{With the board unpowered, use continuity mode to find which two pins are the coil (should read the coil's rated resistance, e.g. 400$\Omega$) and which pins are the switched contacts (should read open/closed as you manually press the relay armature)}{}
\checkitem{Decide: desolder individual relays for a fresh custom board, or keep them in place and reverse-engineer/re-drive the existing traces --- for a first project, desoldering 4--6 relays onto a new small board is far more reliable}{}

% ══════════════════════════════════════════════════════════════════════════════
\roadmapSubSection{1.2}{Desoldering and Testing the Relays}{2--3 Days}

\roadmapHead{Tools}
\begin{toolbadges}
\toolbadge{Soldering iron / hot air}
\toolbadge{Desoldering pump}
\toolbadge{Multimeter}
\toolbadge{Small prototyping/perfboard}
\end{toolbadges}

\checkitem{Heat each of the relay's through-hole pins from the underside while gently pulling from the top --- do not force it, reheat instead of levering}{}
\checkitem{After removal, re-test each relay's coil resistance and contact continuity exactly as in 1.1 --- discard any relay that doesn't click audibly when 12V/5V is briefly applied to the coil}{}
\checkitem{Keep 4--6 confirmed-good relays for this build; save the rest as spares for future projects}{}

\newpage
% ══════════════════════════════════════════════════════════════════════════════
\roadmapSubSection{1.3}{Driver Circuit --- Why You Can't Wire a Relay Straight to a GPIO}{2 Days}

{\small\itshape\color{arligray} This is the single most important lesson of this track: a microcontroller pin cannot directly drive a relay coil.}

\checkitem{An ESP32 GPIO can safely source only a few mA; a 5V/12V relay coil typically wants 30--80mA --- direct connection will not switch it and may damage the pin}{}
\checkitem{Solution: use each GPIO to drive the base/gate of a transistor or MOSFET, which then switches the relay coil's own power supply}{}
\checkitem{Simplest correct approach: a ULN2803A or ULN2003A Darlington-array IC --- 7--8 channels in one chip, built-in flyback diodes, directly accepts 3.3V logic}{}
\checkitem{Wiring per channel: GPIO $\to$ ULN2803 input pin; ULN2803 output pin $\to$ relay coil $-$; relay coil $+$ $\to$ 5V/12V supply; ULN2803 GND $\to$ common ground with ESP32}{}
\checkitem{If building discrete instead of using ULN2803: NPN transistor (e.g. 2N2222/BC547) base through a 1k$\Omega$ resistor from GPIO, collector to relay coil $-$, emitter to GND, and a flyback diode (1N4007) across the coil, cathode to the supply side}{The flyback diode is not optional --- switching off an inductive coil without one generates a voltage spike that can destroy the transistor or the GPIO}

\roadmapHead{Tools to Add}
\begin{toolbadges}
\toolbadge{ULN2803A or ULN2003A IC}
\toolbadge{1N4007 diodes (if going discrete)}
\toolbadge{NPN transistors 2N2222/BC547 (if going discrete)}
\toolbadge{1k$\Omega$ resistors}
\toolbadge{5V/12V bench power supply or salvaged phone charger}
\end{toolbadges}

% ══════════════════════════════════════════════════════════════════════════════
\roadmapSubSection{1.4}{Microcontroller Integration}{2 Days}

\checkitem{Assign one ESP32 GPIO per relay channel (e.g. GPIO 16, 17, 18, 19 for a 4-relay build) --- avoid boot-strapping pins (GPIO 0, 2, 15) for outputs that must stay LOW at boot}{}
\checkitem{Confirm with the multimeter in voltage mode that each GPIO reads ~3.3V HIGH / 0V LOW before wiring to the ULN2803}{}
\checkitem{Wire everything on a breadboard first, verify every channel switches correctly, then move to a soldered perfboard only once proven}{}
\checkitem{Keep the ESP32/logic side and the mains-switched side of the final enclosure physically separated --- this matters even more once real 220V loads are involved (see 1.6)}{}

\newpage
% ══════════════════════════════════════════════════════════════════════════════
\roadmapSubSection{1.5}{Firmware and Control Interface}{3--4 Days}

\roadmapHead{Minimal Firmware --- Local Web Control Panel}
\checkitem{Use the ESP32's built-in WiFi to host a small web server (Arduino \texttt{WebServer.h} or \texttt{ESPAsyncWebServer} library)}{}
\checkitem{Serve a simple HTML page with one button per relay; each button hits a URL like \texttt{/relay1/on} and \texttt{/relay1/off}}{}
\checkitem{In the handler for each route, call \texttt{digitalWrite()} on the matching GPIO and confirm the change back to the page}{}
\checkitem{Reserve a static IP for the ESP32 (see Track 0.4) so the control page is always at the same address, e.g. \texttt{http://192.168.1.50}}{}
\checkitem{Optional upgrade: add a simple password/token check in the web handler before any relay toggles, since this now controls real household devices}{}

\roadmapHead{Optional: Physical Button Backup}
\checkitem{Wire one physical push-button per relay in parallel with the software control, so the panel still works if WiFi is down}{Read the button with \texttt{digitalRead()} and a pull-up resistor (internal \texttt{INPUT\_PULLUP} works on ESP32), toggle the relay state on each press}

% ══════════════════════════════════════════════════════════════════════════════
\roadmapSubSection{1.6}{Final Build --- Enclosure and Mains Safety}{2--3 Days}

\roadmapHead{\textbf{\color{arlizgold}$\bigtriangleup$ CAUTION --- This Step Involves 220V AC}}
\checkitem{Only proceed to switching a real 220V lamp/appliance once every relay has been fully tested at low voltage first}{}
\checkitem{Use a fully enclosed, non-conductive plastic project box; mains wiring must be fully insulated with no exposed copper anywhere inside}{}
\checkitem{Relay contact side wiring: appliance $\to$ relay common (COM) $\to$ relay normally-open (NO) $\to$ mains live, with mains neutral wired straight through --- never switch neutral only}{}
\checkitem{Add an inline fuse sized to the relay's contact rating on the mains side}{}
\checkitem{Keep at least 6--8mm physical separation (creepage distance) between the low-voltage logic section and any 220V wiring inside the enclosure}{}
\checkitem{Before final closing, test with a lamp: toggle each relay from the web panel and the physical button, confirm correct channel mapping and clean on/off switching}{}
\checkitem{Label every channel on the outside of the enclosure to avoid future mix-ups}{}

\vspace{6pt}
\noindent
\begin{tcolorbox}[
  enhanced, colback=bgsubtle, colframe=arlizblue!30,
  leftrule=4pt, sharp corners, left=10pt, right=10pt, top=6pt, bottom=6pt
]
{\small\color{arligray}\itshape
\textbf{\color{arliznavy}Milestone:} A working relay panel, controllable from your phone's browser on the home WiFi,
switching at least one real lamp safely. This proves your solder joints, driver circuit, firmware, and network
setup all work together --- every later track reuses this exact pattern (GPIO $\to$ driver $\to$ real-world load).
}
\end{tcolorbox}
